\section{Eleven integers by intersective interpolation and parity}
\label{sec:eleven-intersective-budget}

The centered-budget guard can also replace the nearest-square guard in
the late twelve-integer intersective construction. A shared Pell
coordinate, an early main-root projection, and four odd radicands then
remove all variable separating weights. The result is an ordinary
polynomial in \textbf{eleven unrestricted integers, of exact degree
$2240$, evaluated in $206$ live binary operations}. Its less-centered
companions have degree $2272$ in $204$ operations and degree $2912$ in
$200$ operations. The late intersective predecessor has twelve integers,
degree $4832$, and $248$ operations. The twelve-integer degree-$1000$
construction remains a lower-degree, lower-operation alternative.

The intersective quintic, full-residue interpolation, and trace-before-
separation have explicit prior sources in the supplied provenance; no
novelty for those individual devices is claimed. The new composition
uses the centered budget, an even shared strong ordinate, fixed
$2$-adic separation, and integral norm centers. All comparisons concern
these concrete equations, not global minimality or a new formal
kernel certificate.

\subsection{Complete definition}

Supply exactly
\[
 j,C,z,h,r,v,w,T,c,n,s\in\mathbb Z.
\]
Use the same six authenticated source34/cone numerals
$\beta,K_0,M_C,M_F,A_c,B_c$, with
$\beta=2^{1064921}-1$, $A_c=2^{1113}$, $B_c=2^{666104}$, and
$M_F$ already shifted. All remaining symbols are polynomial
abbreviations, not additional coordinates:
\begin{align*}
 J&=6j+1,\quad  q=\beta J+1,\quad  W=12w+4,\quad X=qW,\\
 Y&=(2v+1)q^3,\quad  E=XY,\quad a=(X+1)Y,\quad H=4a+3,\\
 \Delta&=a^2+H,\quad F=(K_0+W)C-z(q-1),\quad \alpha=q-F-C,\\
 R&=(q(q-F)-C)(q^2-1)+(M_C+qM_F)J,\\
 k&=R+1+hE,\quad \eta=c-Yk,\quad L=XY^2k,\\
 D_{\rm err}&=(A_cF-q)^2+(B_cC-q)^2,\\
 \mathcal P&=k^2(q^2-D_{\rm err})-4q^2\eta^2,\\
 D_g&=1+4(3\mathcal P-4)s^2,\quad  I=2src^2,\\
 D_f&=1+\Delta I^2,\quad Q=\Delta^2I^2,\quad \delta=Q-1,\\
 V&=D_f(c-R+4IT)-c,\quad C_y=QV^2-1,\quad D_a=\delta C_y,\\
 D_t&=1+L(L+k),\quad D_m=1+\Delta c^2,\quad b_0=X+ac.
\end{align*}
The subscript $a$ in $D_a$ denotes the auxiliary radical, not a new
variable. Set $u=HT-b_0$ and define the two quadratic-algebra coefficients
\begin{align*}
 A_\pi={}&u^5+10u^3D_m+5uD_m^2
       +32H^2(u^3+3uD_m)\\
       &\quad+H^3(u^2+D_m)+31H^4u+31H^5,\\
 B_\pi={}&5u^4+10u^2D_m+D_m^2
       +32H^2(3u^2+D_m)+2H^3u+31H^4.
\end{align*}
In the formal algebra $d^2=D_m$ they satisfy
\[
 A_\pi+B_\pi d
 =((u+d)^2+31H^2)((u+d)^3+H^2(u+d)+H^3).
\]
Finally put
\[
 S=IH^5,\qquad \Theta=(2Q-1)V+4L,\qquad
 Z=S(n+\Theta)-A_\pi.
\]
For the integer polynomial $F_{32}$ defined by its $32$ formal signed
linear factors, the equation is
\begin{equation}
\label{eq:eleven-parity-norm}
 F_{32}\bigl(Z;
 B_\pi^2D_m,\,4S^2D_a,\,S^2D_f,\,16S^2D_t,\,64S^2D_g\bigr)=0.
\end{equation}
Here every radicand is an integer polynomial and there is no denominator
to clear. The literal circuit uses only binary addition, subtraction,
and multiplication. The direct and auxiliary-only companions replace
$\Theta$ by $0$ and $(2Q-1)V$, respectively. These are invertible
integer-polynomial translations in $n$, not additional constraints.

\subsection{The intersective polynomial and the zero-scale exclusion}

The prior interpolation construction uses
\[
 \pi(x)=(x^2+31)(x^3+x+1)
       =x^5+32x^3+x^2+31x+31.
\]
It has a root modulo every nonzero integer, no zero in a real extension
of $\mathbb Q$ of degree at most two, and positive derivative on
$\mathbb R$. For completeness, at an odd prime other than $31$, either
the cubic has a root or its irreducible Frobenius $3$-cycle makes its
discriminant $-31$ a square. These roots are simple. At $31$, the cubic
has simple root $3$; at $2$, $-31\equiv1\pmod8$ is a $2$-adic square.
Prime-power lifting and CRT give intersectivity. The quadratic has no
real root and the cubic has no rational root, so neither factor has a
zero in a real field of degree at most two. Finally
\[
 \pi'(x)=5x^4+96x^2+2x+31>0.
\]
For $H\ne0$, write $\Phi_H(x)=H^5\pi(x/H)$. Then
$A_\pi\pm B_\pi\sqrt{D_m}=\Phi_H(u\pm\sqrt{D_m})$ are nonzero.
Moreover $B_\pi>0$, either by monotonicity or by the exact certificate
\[
 B_\pi=5u^4+10u^2D_m+D_m^2+32H^2D_m
       +95H^2u^2+H^2(u+H)^2+30H^4>0.
\]

For every integer input, $q$ is even, $|q|\ge16$, $|X|\ge64$,
$|Y|\ge4096$, and $|a|\ge63\cdot4096$. Thus $H\ne0$,
$\Delta>1$ is odd, and $D_m,D_f,D_t$ are positive. Also
$Q=(\Delta I)^2\ne1$, so $\delta\ne0$. If $I=0$, $D_a=1$.
If $I\ne0$, then $c\ne0$, $|I|\ge c^2$ and $D_f>|c|+1$.
For the integer $b=c-R+4IT$, $V=D_fb-c$ has $|V|\ge1$.
Therefore $C_y>0$ and $D_a>0$.

At a zero of \eqref{eq:eleven-parity-norm}, $S=0$ would specialize the
norm to $(A_\pi^2-B_\pi^2D_m)^{16}\ne0$. Hence $S\ne0$, and
$I,s,r,c$ are nonzero. If $D_g<0$, a vanishing complex factor would
have precisely one nonzero imaginary summand, with coefficient $8S$.
This is impossible. Thus $D_g\ge0$, which for nonzero integer $s$
forces $\mathcal P\ge2$. In particular
\[
 D_{\rm err}<q^2,\qquad |\eta/k|<\tfrac12.
\]
The disk inequality implies $0<A_cF/q<2$ and $0<B_cC/q<2$.
Since $A_c,B_c\ge8$, $F,C,\alpha$ all have the same nonzero sign as
$q$. Also
\[
 |k|(|Y|-\tfrac12)<|c|<|k|(|Y|+\tfrac12),
 \qquad \operatorname{sign}(c)=\operatorname{sign}(kY).
\]
Positivity of $\eta$ has not been presumed.

\subsection{Early main-root projection and parity extraction}

A vanishing factor can be written, with independently signed roots, as
\[
 S(n+\Theta)=A_\pi+B_\pi d
             +S(2\widetilde y+f+4\tau+8g).
\]
All five underlying radicands are nonnegative integers. Project this
identity onto the rational squareclass. Every rational radical is an
integer, while irrational squareclass terms have zero rational
projection. If $d=\pm\sqrt{D_m}$ were irrational, it would follow that
$S\mid A_\pi$, hence $H\mid A_\pi$.

But $q\equiv2$, $W\equiv1$, and $X\equiv2\pmod3$, so
\[
 a\equiv H\equiv\Delta\equiv0,\qquad
 b_0\equiv2,\quad u\equiv D_m\equiv1\pmod3.
\]
Consequently $A_\pi\equiv1+1+2\equiv1\pmod3$, contradicting
$H\mid A_\pi$. Thus the actual signed main root $d$ is integral.
Projecting again yields
\[
 IH^5\mid A_\pi+B_\pi d=\Phi_H(d-b_0+HT).
\]
Monicity forces $H\mid d-b_0$: if $p^e\Vert H$ and the argument had
valuation $b<e$, its leading fifth power would have valuation $5b$,
strictly lower than every other term, contradicting divisibility by
$p^{5e}$. Hence $\gamma=(d-b_0)/H$ is integral.

Now subtract the main-root contribution and divide by $S$. The signed
sum $f+2\widetilde y+4\tau+8g$ is rational. All four radicands are odd:
\[
 D_f\equiv D_a\equiv D_g\equiv1\pmod4,
 \qquad D_t\equiv1\pmod8.
\]
Indeed $I$ is even, $Q\equiv0\pmod4$ and
$\delta\equiv C_y\equiv-1\pmod4$; $X$ is a multiple of $8$.
A positive odd integer radicand is an odd square times an odd squarefree
integer. Within any nontrivial squareclass, the participating signed
coefficients therefore have distinct $2$-adic valuations from
$0,1,2,3$. The unique least valuation cannot cancel. Independence of
distinct squarefree radicals implies that all four roots are rational,
and hence integers. No comparison of their magnitudes is required.

The rational $y=\widetilde y/\delta$ satisfies
$QV^2-(Q-1)y^2=1$ and is integral by adjacent-coefficient descent.
Explicitly, a reduced denominator $b$ has $b^2\mid Q-1$; the
norm-preserving transformation
\[
 V'=(2Q-1)V-2(Q-1)y,\qquad
 y'=(2Q-1)y-2QV
\]
keeps $V'$ integral and the denominator unchanged, while reducing
$|V|$ to $1$. The terminal value has $|y|=1$, proving integrality.

\subsection{Soundness with an even strong multiplier}

The recovered equations and the exact definitions give
\begin{gather*}
 f^2-\Delta I^2=1,\qquad
 \rho^2=\Delta(f^2-1),\qquad
 \rho^2V^2-(\rho^2-1)y^2=1,\qquad \rho=\Delta I,\\
 d=X+ac+H\gamma,\quad d^2-\Delta c^2=1,\quad
 \tau^2-XY^2(XY^2+1)k^2=1,\\
 V=fo-c=cj_{\rm aux}-R,\qquad k=R+1+hE,
\end{gather*}
where
\[
 o=f(c-R+4IT),\qquad
 j_{\rm aux}=4\Delta s^2r^2c^3(c-R)+8srcD_fT.
\]
The incoming multiplier is even. We do not assert a nonintegral
restoration such as $(2\Delta sr-1)/2$.

Here is the exact modification of the centered source's signed rank
proof. Its common-sign packing bounds give, for $Q_0=|q|$,
\[
 Q_0(Q_0^2-1)<R<Q_0^4-2,
 \qquad |E|>3R+1.
\]
Put $B_0=|a+2|$ and $P_0=|2XY^2+1|$. First/main Pell classification,
using only the already proved centered interval and sign identity,
gives $|k|=2\psi_{P_0}(n_0)$, $|c|=\psi_{B_0}(p)$,
$p>n_0\ge24$, and $|c|>2p,2R$. The first-index congruence has
positive branch $n_0\ge(R+1)/2$ and negative branch $n_0>R$.

The strong norm gives $|I|=\psi_{B_0}(m)$ and
$|f|=\chi_{B_0}(m)$. Since $c^2\mid I$, strong divisibility first
implies $p\mid m$. Write $m=pt$ and $D=\chi_{B_0}(p)$. Composition
and binomial expansion give
\[
 |I|=|c|\psi_D(t),\quad D^2-1=\Delta c^2,\quad
 \psi_D(t)\equiv tD^{t-1}\pmod{|c|}.
\]
Since $D$ is coprime to $c$, $|c|\mid t$. Thus $p|c|\mid m$ and
$m>2p$. These are precisely the strong-rank inputs; oddness was not
used.

For the auxiliary parameter $|\rho|>1$, its Pell index $\ell$ is odd
because $\chi_{|\rho|}(2t)\equiv(-1)^t\pmod{|\rho|}$, while its
abscissa is divisible by $|\rho|$. The exact odd-index quotient
identities, together with $V\equiv-c\pmod{|f|}$ and
$V\equiv-R\pmod{|c|}$, yield
\[
 \psi_{B_0}(\ell)\equiv\pm\psi_{B_0}(p)\pmod{|f|},
 \qquad \ell\equiv\pm R\pmod{|c|}.
\]
The signed abscissa step-down then gives
$\ell\equiv\pm p\pmod m$. Since $|c|\mid m$ and
$0<p,R<|c|/2$, $p=R$. The negative first-index branch is impossible;
no-wrap gives $2n_0=R+1$. The unchanged sign-uniform ratio and main-root
congruence prove $|Y|>2^R$ and $X=\operatorname{sign}(c)2^R$.
Since $X\equiv2\pmod3$ and $R$ is odd, $X=2^R$, $c>0$;
the surviving first-index sign then gives $k,Y>0$.

The centered source's uniform half-binomial estimate now applies.
Writing $R=2t+1$, it gives an even integer $M_{\rm bin}$ with
\[
 \frac{M_{\rm bin}}2<\frac ck<\frac{M_{\rm bin}}2+\frac14.
\]
Its proof bounds the reciprocal binomial tail by $1/8$ and the Pell
ratio error by $12t/(2^R+1)<1/8$, uniformly for the prescribed $Y$.
Since $|c/k-Y|<1/2$ and $Y$ is integral,
$Y=M_{\rm bin}/2$ and $0<\eta/k<1/4$. Thus $\eta,k-\eta>0$.

An exact old odd-coordinate chart can now be restored with fresh
private witnesses, rather than by changing the parity of incoming $I$.
Let $A_0=a+2$ and put
$I_* =\psi_{A_0}(Rc)$, $f_* =\chi_{A_0}(Rc)$.
Here $A_0$ is even, $R,c$ are odd, and composition gives $c^2\mid I_*$;
therefore $I_*$ and $\rho_* =\Delta I_*$ are odd. If $t=(R-1)/2$ and
$Q_t(z)=\chi_{\sqrt z}(2t+1)/\sqrt z$ is the odd-index quotient
polynomial, choose
\[
 V_*=(-1)^{t+1}Q_t(\rho_*^2),\qquad y_*=\psi_{\rho_*}(R).
\]
Its two evaluations give $V_*\equiv-c\pmod{f_*^2}$ and
$V_*\equiv-R\pmod c$. Hence the old strong and auxiliary coordinates
are integral, with odd multiplier $\rho_*/c^2$, and preserve
$k=R+1+hE$. The recovered positive guards and $X=2^R$ restore the
remaining old rows. The pinned ordinary native/compiler and
finite-refutation theorem therefore applies. This regeneration works
for every odd $R$ on soundness, not only $R\equiv3\pmod4$.

\subsection{Full-residue interpolation and canonical completeness}

We use the prior full-residue lemma in its precise order: for even
$A_0>1$, $R\equiv3\pmod4$, $c=\psi_{A_0}(R)\equiv R\pmod4$, and
any prescribed $N>0$, one can first fix $I,f>0$ with
\[
 f^2-(A_0^2-1)I^2=1,\qquad 4Nc^2\mid I,
\]
such that \emph{every} residue class modulo this same $I$ contains a
$T$ for which
\[
 V=f^2(c-R+4IT)-c,\qquad
 ((A_0^2-1)I)^2V^2-
 \bigl(((A_0^2-1)I)^2-1\bigr)y^2=1
\]
has an integer solution.

For precision, finite order of the Pell recurrence modulo $4Nc^2$
selects a positive even index $m_0$ and such a fixed $I,f$.
The correct period modulo $f^2$ is $\lambda=4m_0f$, not merely
$4m_0$. Indeed, with $\alpha=A_0+\sqrt{A_0^2-1}$,
$\alpha^{4m_0}\equiv1-4f\alpha^{m_0}\pmod{f^2}$, whose $f$th power
is $1$. Lucas lifting gives $v_p(I)\ge v_p(m_0)$ for $p\mid I$;
at $2$ the exact formula is $v_2(I)=v_2(m_0)+v_2(A_0)$.
Thus $\gcd(\lambda,4I^2)\mid4I$.

For a requested class $T\equiv a\pmod I$, the congruences
\[
 \ell\equiv R\pmod\lambda,\qquad
 \ell\equiv R-4Ia\pmod{4I^2}
\]
are compatible. Choose a positive solution, which is $3$ modulo $4$.
For $S_0=(A_0^2-1)I$, put
$V_0=\chi_{S_0}(\ell)/S_0$ and $y=\psi_{S_0}(\ell)$.
The odd quotient identities give
$V_0\equiv-c\pmod{f^2}$ and
$V_0\equiv-\ell\pmod{4S_0^2}$, including the important factor $4$.
Consequently
\[
 T=\frac{V_0+c-f^2(c-R)}{4If^2}
\]
is integral and is $a$ modulo $I$, since $f^2\equiv1\pmod I$.
This proves the simultaneous fixed-strong-witness interpolation needed
below; its full valuation proof is retained locally with the precursor.

Start now with the authentic canonical source34 family. Its prescribed
half-binomial $Y$ gives
$D_{\rm err}\le q^2/2$ and $0<\eta/k<1/4$, so
$\mathcal P>q^2k^2/4\ge64$. Put $D=3\mathcal P-4>0$.
Since $D\equiv2\pmod3$, it is nonsquare. Square any positive Pell
unit $g_0+u_0\sqrt D$ to obtain
\[
 g=g_0^2+Du_0^2,\qquad s=g_0u_0>0,
 \qquad g^2-4Ds^2=1.
\]
Apply the interpolation lemma with $N=2s$; it fixes $I,f$ with
$8sc^2\mid I$. Set $r=I/(2sc^2)$. Choose an intersective root
$a$ modulo $I$ of $\pi(a+\gamma)$; the already fixed $I,f$ then
supply $T\equiv a\pmod I$ and the auxiliary root $y$.
With $\widetilde y=\delta y$, all five radicands are integer squares,
and
\[
 A_\pi+B_\pi d=H^5\pi(T+\gamma)
\]
is divisible by $S$. Finally set
\[
 n=\frac{A_\pi+B_\pi d}{S}
        +2\widetilde y+f+4\tau+8g-\Theta.
\]
This integer makes a norm factor vanish. The witness order is native
data and budget, then $s$, then fixed $I,f,r$, then the intersective
class and $T,y$, then $n$. Nothing chosen later changes the positive
budget or reselects the strong witness after the desired class is known.

\subsection{Exact metrics and the inherited conclusion}

\begin{center}
\begin{tabular}{lrrr}
Center in the $n$ carrier & Integers & Exact degree & Binary operations\\
\hline
$0$ & $11$ & $2912$ & $200$\\
$(2Q-1)V$ & $11$ & $2272$ & $204$\\
$(2Q-1)V+4L$ & $11$ & $2240$ & $206$
\end{tabular}
\end{center}

The degree certificates are independent of the existence proof.
The exact identities behind the centers include
\[
 4D_a-((2Q-1)V)^2=-V^2-4Q+4,\qquad
 16D_t-(4L)^2=16Lk+16.
\]
The supplied leading-form calculation gives degree $2240$ for the
primary equation. On the actual fixed-numeral slice it has the
nonzero coefficient
\begin{align*}
 &[j^{1176}w^{280}v^{296}h^{16}s^{120}r^{112}c^{224}T^{16}]\,P_{11}\\
 &\hspace{2em}=2^{2592}3^{1460}\beta^{1176}>0.
\end{align*}
The exponent sum is $2240$. The exact sparse degree cut agrees with the
formal root filtration $16\cdot91+8\cdot51+8\cdot47=2240$;
the supplied coefficient proves that this upper bound is attained.
The literal circuits and independent checkers
report every live binary operation, with no uncounted division or root.

Under the same explicitly inherited ordinary native/compiler and
calibration theorems, ZFC proves
\[
 [\forall\mathbf x\in\mathbb Z^{11},\ P_{11}(\mathbf x)\ne0]
 \quad\Longleftrightarrow\quad\operatorname{Con}(\mathrm{ZFC}).
\]
The result neither assumes nor proves consistency. The unbounded
reduction, the inherited calibration, finite circuit and coefficient
checks, and the separate proof-assistant/object-ZFC obligations remain
distinct. The complete local source pins and retained notices identify
which parts of the interpolation, centered-budget, and signed-Pell
arguments are inherited.
